% =====================================================================
%  M5 -- GD against Newton
% =====================================================================
\subsection{M5: GD against Newton}
\setprob{M5}

\begin{frame}{\probtitle{M5}{GD against Newton}}
  \begin{probbox}[Problem M5 \hfill Mathematical \quad $\star\star\star$]
    \begin{enumerate}[(a)]
      \item Apply one Newton step, $x_{k+1} = x_k - f'(x_k)/f''(x_k)$, to $f(x) = (x-3)^2$ from $x_0 = 0$.
            How many GD steps with $\alpha = 0.1$ reach the same accuracy (error cut by $10^6$)?
      \item Now take $J(x) = x^4 - 6x^3 + 8x^2$ from Problem M1 and start at $x_0 = 1.4$.
            Run Newton and GD ($\alpha = 0.05$). Where does each one go?
      \item Which method would you trust here, and why?
    \end{enumerate}
  \end{probbox}
  \footnotesize\textcolor{mGrey}{Practice: comparing GD and Newton with numerical examples.}
\end{frame}

\begin{frame}{\soltitle{M5}{1}{3}}
  \begin{columns}[T]
    \begin{column}{0.52\textwidth}
      \small
      \textbf{(a) Newton on a parabola.} $f'(0) = -6$ and $f'' = 2$:
      \[
        x_1 = 0 - \frac{-6}{2} = \ans{3}
      \]
      Newton reaches the minimum. Newton fits a parabola to $f$, and here $f$ \emph{is} a parabola.

      \vspace{3pt}
      GD with $\alpha = 0.1$ cuts the error by $0.8$ per step, so it needs $\ans{62}$ steps
      (Problem M2). Newton's step is $\alpha = 1/f''$, the best possible fixed rate.
    \end{column}
    \begin{column}{0.46\textwidth}
      \centering
      \includegraphics[width=\linewidth, height=0.6\textheight, keepaspectratio]{fig_gd_vs_newton}
    \end{column}
  \end{columns}
\end{frame}

\begin{frame}{\soltitle{M5}{2}{3}}
  \textbf{(b) Start at $x_0 = 1.4$ on the quartic.} Here $J'(1.4) = -1.904$ and $J''(1.4) = -10.88 < 0$:
  the curve bends \emph{down}.

  \vspace{3pt}
  \textbf{Newton:} $x_1 = 1.4 - \dfrac{-1.904}{-10.88} = 1.4 - 0.175 = 1.225$, then
  \[
    1.4 \to 1.225 \to 1.21923 \to \ans{1.21922} \quad\text{the \textcolor{mRed}{maximum}, } J = 3.227 .
  \]
  \textbf{GD:} $x_1 = 1.4 - 0.05(-1.904) = 1.4952$, then
  \[
    1.4 \to 1.4952 \to 1.6426 \to 1.8704 \to\cdots\to \ans{3.2808}\quad\text{the global minimum, } J = -9.915 .
  \]
  Newton converged quickly to a maximum.
\end{frame}

\begin{frame}{\soltitle{M5}{3}{3}}
  \textbf{(c) Why.} Newton solves $J'(x) = 0$. It cannot tell a minimum from a maximum, and where $J'' < 0$
  its step points \emph{uphill}. GD steps along $-J'$, which always points downhill, so $J$ decreases for any
  small enough $\alpha$.
  \begin{center}\small
  \begin{tabular}{@{}l c c@{}}
    \toprule
    from $x_0 = 1.4$ & Newton & GD, $\alpha = 0.05$ \\
    \midrule
    limit & $1.2192$ (maximum) & $3.2808$ (minimum) \\
    speed near the limit & quadratic & linear \\
    \bottomrule
  \end{tabular}
  \end{center}
  \begin{keybox}[Take-away]
    Trust GD far from the answer, and Newton close to a minimum where $J'' > 0$.
    Levenberg--Marquardt~\cite{nocedal2006} uses both gradients and curvature.
  \end{keybox}
\end{frame}
